\begin{explanationbox}
\textbf{\textcolor{CExplanation}{一句话直觉}}

\textbf{椭圆是把圆在横、纵方向分别拉伸得到的；圆的切线“半径垂直切线”，拉伸后就变成椭圆的切线公式。}

\medskip
\textbf{\textcolor{CExplanation}{第一层：为什么必须先确认 $P$ 在椭圆上？}}

我们讨论的是“\textbf{在 $P$ 处}的切线”。这首先要求 $P$ 是椭圆上的点：
\[
\frac{x_0^2}{a^2}+\frac{y_0^2}{b^2}=1.
\]
若点在椭圆外，可从该点作两条切线，但外点本身不是切点；直接将外点坐标代入切点公式，得到的不是“在外点处的切线”。

\medskip
\textbf{\textcolor{CExplanation}{第二层：先回到最熟悉的圆}}

单位圆 $u^2+v^2=1$ 上一点 $P'(u_0,v_0)$，半径 $OP'$ 与切线垂直。
故 $P'$ 处切线上任意点 $Q'(u,v)$ 满足
\[
(u-u_0)u_0+(v-v_0)v_0=0.
\]
因为 $u_0^2+v_0^2=1$，所以
\[
\boxed{u_0u+v_0v=1.}
\]
这正是\textbf{圆的切线方程}，没有使用微积分。

\medskip
\textbf{\textcolor{CExplanation}{第三层：把椭圆压回圆（数形结合）}}

对椭圆引入新坐标
\[
u=\frac{x}{a},\quad v=\frac{y}{b},\quad
u_0=\frac{x_0}{a},\quad v_0=\frac{y_0}{b}.
\]
原椭圆立即变成 $u^2+v^2=1$，切线 $u_0u+v_0v=1$ 还原为
\[
\frac{x_0}{a}\frac{x}{a}+\frac{y_0}{b}\frac{y}{b}=1
\quad\Longrightarrow\quad
\boxed{\frac{x_0x}{a^2}+\frac{y_0y}{b^2}=1.}
\]
\textbf{形象记忆：} 不是硬背平方替换，而是\textbf{先把椭圆看成圆，再用圆的切线公式}。

\medskip
\textbf{\textcolor{CExplanation}{第四层：切线究竟“垂直”于哪个方向？}}

椭圆切线的一组法向量（与切线垂直的向量）是
\[
\boldsymbol n=\left(\frac{x_0}{a^2},\frac{y_0}{b^2}\right).
\]
因此切线过 $P$ 且满足
\[
\boldsymbol n\cdot\overrightarrow{PQ}=0
\iff
\frac{x_0}{a^2}(x-x_0)+\frac{y_0}{b^2}(y-y_0)=0.
\]
把切点满足的椭圆方程移到右边，就得到核心公式。

\textbf{注意：} 对\textbf{圆}，法向量 $(x_0,y_0)$ 就是半径方向；但对\textbf{非圆椭圆}，一般不能把 $OP$ 当作法线。纵横两个半轴的缩放比例不同，导致法向量多出 $a^2,b^2$ 两个分母。

\medskip
\textbf{\textcolor{CExplanation}{第五层：为什么切线只碰到椭圆一次？}}

在椭圆上，令两个单位向量
\[
\boldsymbol p=\left(\frac{x_0}{a},\frac{y_0}{b}\right),\quad
\boldsymbol q=\left(\frac{x}{a},\frac{y}{b}\right).
\]
柯西不等式告诉我们 $\boldsymbol p\cdot\boldsymbol q\le1$，且等号只在 $\boldsymbol p=\boldsymbol q$ 时成立。
这就是说：椭圆所有点都在直线
$\dfrac{x_0x}{a^2}+\dfrac{y_0y}{b^2}=1$ 的同一侧，唯独 $P$ 在直线上。
所以它是椭圆的\textbf{支撑直线}，也就是切线。

\medskip
\textbf{\textcolor{CExplanation}{第六层：和其他圆锥曲线的联系}}
\begin{itemize}[leftmargin=*, itemsep=0.3em]
 \item 圆 $x^2+y^2=r^2$：切线 $x_0x+y_0y=r^2$。
 \item 椭圆 $x^2/a^2+y^2/b^2=1$：切线 $x_0x/a^2+y_0y/b^2=1$。
 \item 双曲线 $x^2/a^2-y^2/b^2=1$：切线 $x_0x/a^2-y_0y/b^2=1$。
 \item 抛物线 $y^2=2px$：切线 $y_0y=p(x+x_0)$（$p\ne0$，切点在曲线上）。
\end{itemize}
\textbf{提醒：} 抛物线的线性项 $2px$ 需要同时考虑 $x$ 与 $x_0$，不能简单照搬“仅替换平方项”。

\medskip
\textbf{\textcolor{CExplanation}{解题判断顺序}}

\textbf{已知切点} $\Rightarrow$ 直接套切点公式；\quad
\textbf{已知斜率} $\Rightarrow$ 用 $m^2=a^2k^2+b^2$；\quad
\textbf{已知椭圆外一点} $\Rightarrow$ 先求切点或设过外点的切线。
\end{explanationbox}
